%% SYNTHETIC TEST UNIT generated by writing/tools/make_outline_book.py from FINAL_OUTLINE_v2 -- not book text
\chapter{因子化的数学与场论工具}\label{ch:prelim}
本单元为测试用的合成文本（由 FINAL\_OUTLINE\_v2 自动生成），陈述内容从略。

\section{第 1 节（测试）}\label{sec:prelim-lightfront}

\begin{definition}[光锥坐标与 Sudakov 分解]\label{def:prelim-lightfront}
（测试用合成陈述：内容从略。）
% uses: none
\end{definition}

\begin{lemma}[推动下的标度 \tagR]\label{lem:prelim-boost}
（$R_{\mathrm{def}}$：由算符定义直接推出。）
（测试用合成陈述：内容从略。）
\uses{def:prelim-lightfront}
\end{lemma}

\begin{example}[共线动量与好/坏分量的标度]\label{ex:prelim-boost-scaling}
（测试用合成陈述：内容从略。）
\uses{lem:prelim-boost}
\end{example}

\begin{definition}[Wilson 线]\label{def:prelim-wilson}
（测试用合成陈述：内容从略。）
\uses{def:prelim-qcd, def:prelim-lightfront}
\end{definition}

\begin{proposition}[Wilson 线的规范协变性 \tagR]\label{prop:prelim-wilson-covariance}
（$R_{\mathrm{def}}$：由算符定义直接推出。）
（测试用合成陈述：内容从略。）
\uses{def:prelim-wilson, prop:qcd-covariance}
\end{proposition}

\section{第 2 节（测试）}\label{sec:prelim-wilson}

\begin{proposition}[半无限直 Wilson 线的动量空间 Feynman 规则 \tagR]\label{prop:prelim-wilson-rules}
（$R_{\mathrm{def}}$：由算符定义直接推出。）
（测试用合成陈述：内容从略。）
\uses{def:prelim-wilson, def:prelim-feynman-rules, prop:appB-rules}
\end{proposition}

\begin{lemma}[程函恒等式]\label{lem:prelim-eikonal-identity}
（测试用合成陈述：内容从略。）
\uses{lem:appB-shuffle}
\end{lemma}

\begin{example}[$m=2$ 的程函恒等式与入射夸克传播子的程函极限]\label{ex:prelim-eikonal-two}
（测试用合成陈述：内容从略。）
\uses{lem:prelim-eikonal-identity, prop:prelim-wilson-rules}
\end{example}

\begin{example}[有限类光线段的单胶子项]\label{ex:prelim-wilson-segment}
（测试用合成陈述：内容从略。）
\uses{def:prelim-wilson, prop:prelim-wilson-rules}
\end{example}

\begin{definition}[维数正规化的技术约定]\label{def:prelim-dr-conventions}
（测试用合成陈述：内容从略。）
\uses{def:qcd-dimreg, def:prelim-msbar}
\end{definition}

\section{第 3 节（测试）}\label{sec:prelim-dimreg}

\begin{lemma}[无标度积分为零]\label{lem:prelim-scaleless}
（测试用合成陈述：内容从略。）
\uses{def:qcd-dimreg}
\end{lemma}

\begin{example}[无标度积分的紫外/红外记账]\label{ex:prelim-scaleless}
（测试用合成陈述：内容从略。）
\uses{lem:prelim-scaleless, def:prelim-dr-conventions}
\end{example}

\begin{definition}[$\gamma_5$ 与 $\varepsilon$ 张量]\label{def:prelim-gamma5}
（测试用合成陈述：内容从略。）
\uses{def:qcd-dimreg}
\end{definition}

\begin{definition}[加号分布]\label{def:prelim-plus}
（测试用合成陈述：内容从略。）
% uses: none
\end{definition}

\begin{lemma}[$(1-z)^{-1-a\varepsilon}$ 的分布展开]\label{lem:prelim-plus-expansion}
（测试用合成陈述：内容从略。）
\uses{def:prelim-plus}
\end{lemma}

\section{第 4 节（测试）}\label{sec:prelim-distributions}

\begin{example}[单圈实辐射的典型积分]\label{ex:prelim-plus}
（测试用合成陈述：内容从略。）
\uses{lem:prelim-plus-expansion}
\end{example}

\begin{definition}[Mellin 变换与卷积]\label{def:prelim-mellin}
（测试用合成陈述：内容从略。）
\uses{def:prelim-plus}
\end{definition}

\begin{lemma}[Mellin 卷积定理与加号分布的矩]\label{lem:prelim-mellin}
（测试用合成陈述：内容从略。）
\uses{def:prelim-mellin, def:prelim-plus, lem:appA-mellin-table}
\end{lemma}

\begin{remark}[矩问题工具箱]\label{rem:prelim-moment-tools}
（测试用合成陈述：内容从略。）
\uses{lem:appK-hausdorff, thm:appK-even-moments, cor:appK-pdf, thm:appK-carlson, prop:appK-continuation, thm:appK-muntz, rem:appK-counterexamples}
\end{remark}

\begin{definition}[割图]\label{def:prelim-cut}
（测试用合成陈述：内容从略。）
\uses{def:prelim-feynman-rules, prop:prelim-wilson-rules}
\end{definition}

\section{第 5 节（测试）}\label{sec:prelim-moments}

\begin{theorem}[切割规则 \tagR]\label{thm:prelim-cutting-rules}
（$R_{\mathrm{PT}}$：在重整化微扰论中对红外正规化的部分子态严格。）
（测试用合成陈述：内容从略。）
\uses{def:prelim-cut, def:prelim-feynman-rules}
\end{theorem}

\begin{proposition}[物理幺正性与鬼场割线 \tagR]\label{prop:prelim-unitarity}
（$R_{\mathrm{PT}}$：在重整化微扰论中对红外正规化的部分子态严格。）
（测试用合成陈述：内容从略。）
（本命题只涉及部分子态。）
\uses{def:prelim-cut, thm:prelim-cutting-rules, prop:qcd-brst-cohomology, thm:prelim-slavnov-taylor, prop:appJ-ghost-cuts, ex:qcd-qqbar-gg}
\end{proposition}

\begin{remark}[无质量维数正规化中的割图规则 \tagA]\label{rem:prelim-cutting-dr}
（测试用合成陈述：内容从略。）
\uses{thm:prelim-cutting-rules, prop:prelim-unitarity, def:qcd-dimreg}
\end{remark}
\begin{proofsketch}测试。\end{proofsketch}
\paragraph{地位与缺口}测试：本陈述的地位由\cref{prop:ee-cutting-exact}、\cref{asm:appN-cutsum}决定。

\begin{proposition}[流关联函数的光学定理 \tagR]\label{prop:prelim-optical}
（$R_{\mathrm{def}}$：由算符定义直接推出。）
（测试用合成陈述：内容从略。）
非微扰地位：\cref{asm:intro-wightman}。
\uses{def:prelim-cut, asm:intro-wightman}
\end{proposition}

\begin{example}[维数正规化中的单圈无质量流关联函数]\label{ex:prelim-vacpol}
（测试用合成陈述：内容从略。）
\uses{def:prelim-dr-conventions, def:prelim-msbar, thm:prelim-cutting-rules, prop:prelim-optical, lem:appC-bubble}
\end{example}
